\section{Compresión de redes neuronales}
\label{sec:compresion_de_redes}

La compresión de redes neuronales es un proceso que consiste en reducir el tamaño de los datos, en particular, de los parámetros de una red neuronal. Esta compresión busca disminuir la complejidad del modelo en dos dimensiones: por un lado, la complejidad espacial, relacionada con la memoria necesaria para almacenar el modelo y el otro, la complejidad temporal, asociada al costo computacional de las operaciones.

En este sentido, la complejidad espacial se define como la cantidad de memoria necesaria para almacenar los parámetros del modelo. Al reducir el tamaño de cada uno de estos parámetros, se disminuye la cantidad total de memoria requerida. Esto resulta especialmente beneficioso para dispositivos con recursos limitados, como dispositivos móviles o sistemas embebidos, donde la memoria disponible es un recurso crítico \cite{DBLP:journals/corr/abs-2010-03954}.

Por su parte, la complejidad temporal se mide en términos de la cantidad de FLOPs necesarias para realizar una inferencia con el modelo. La reducción de la cantidad de parámetros o de su precisión disminuye la complejidad de los cálculos necesarios y, por ende, la cantidad de operaciones de procesador requeridas. Esto acelera el tiempo de inferencia \cite{molchanov2017pruning}.

La compresión también mejora la eficiencia energética del modelo, ya que el acceso a memoria y la realización de operaciones computacionales son procesos que consumen energía. Al reducir la cantidad de datos que deben ser procesados y almacenados, se disminuye la cantidad de FLOPs necesarias y, por ende, se disminuye el consumo energético asociado a la inferencia del modelo. Esto es beneficioso tanto para sistemas con recursos limitados como para aplicaciones que buscan minimizar su huella energética, como la computación en la nube. En los últimos años compañías como OpenAI, Google, entre otras, han disponibilizado sus modelos en la nube en lo que se conoce como \textit{Machine Learning as a Services} (MLaaS), donde la eficiencia energética es un factor clave para reducir costos operativos y el impacto ambiental \cite{barros2024scheduling}.

\TODO{Revisar este párrafo luego de presentar los resultados, porque creo que vamos a probar con 4, 6 y 8 bits, ¿no?
En ese caso esto no aplica, más allá de que no le demos tanta bola.}
Dado el alcance de esta investigación, este trabajo se enfoca en evaluar el impacto sobre la complejidad espacial, dado que constituye el factor determinante para el despliegue de modelos en dispositivos con memoria limitada. La complejidad temporal, aunque importante, no se aborda en profundidad, ya que su exploración cuantitativa está fuera del alcance del mismo.

\subsection{Técnicas de compresión}
Existen múltiples técnicas de compresión de redes neuronales, las cuales permiten reducir la cantidad de parámetros y operaciones necesarias para la inferencia, sin afectar significativamente la precisión del modelo. Algunas de las técnicas más comunes son:

\begin{itemize}
    \item Compartición de pesos (\textit{Weight Sharing} en inglés): se agrupan los pesos de la red en un conjunto reducido de valores, lo que permite una representación más compacta sin afectar demasiado la capacidad de generalización \cite{8253600}.
    \item Poda (\textit{Weight Pruning} en inglés): se eliminan conexiones o neuronas que tienen una contribución marginal en la inferencia, lo que reduce la cantidad de parámetros sin afectar significativamente la precisión del modelo \cite{8253600}.
    \item Cuantización (\textit{Quantization} en inglés): se reduce la cantidad de bits utilizados para representar los parámetros del modelo, lo que disminuye el consumo de memoria y acelera las operaciones de inferencia \cite{8253600}.
    \item Factorización de tensores (\textit{Tensor Decomposition} en inglés): se aproximan los tensores de pesos mediante técnicas de factorización, lo que reduce la cantidad de FLOPS necesarias durante la inferencia \cite{8253600}.
    \item Destilación de conocimiento (\textit{Knowledge Distillation} en inglés): se entrena un modelo más pequeño (\textit{student}) para imitar el comportamiento de un modelo más grande (\textit{teacher}), lo que transfiere la información esencial en una forma más compacta \cite{8253600}.
\end{itemize}

Cada técnica posee ventajas y desventajas y la elección de la técnica adecuada depende del caso de uso específico, el tipo de arquitectura del modelo y los requisitos de rendimiento \cite{DBLP:journals/corr/abs-2010-03954}.

\subsubsection{Cuantización}
\label{sec:cuantizacion}

La cuantización es una compresión con perdida y tiene su origen en el procesamiento de señales digitales, donde se utiliza para representar señales continuas en un formato discreto. En el dominio analógico, las señales pueden tomar cualquier valor dentro de un rango continuo, mientras que en el dominio digital, estas señales deben ser representadas mediante un conjunto finito de valores discretos, la transformación específica de estos valores depende de la función de cuantización.

La función de cuantización $q(x)$ mapea un valor continuo de entrada $x$ a un conjunto discreto de salida $y$. Esta función queda definida por dos series de parámetros: los niveles de cuantización $\mathbf{l}$, que definen los valores posibles de salida y los umbrales de cuantización $\mathbf{t}$ que delimitan el espacio de entrada en m\'ultiples intervalos y se asigna a cada uno de ellos un nivel de cuantización específico \cite{sayood2017introduction}.

Por ejemplo, podemos definir una función de cuantización de $8$ niveles, con $8$ intervalos como se muestra en la tabla \ref{tab:quantization_example}.

\begin{table}[H]
    \centering
    \caption{Función de cuantización de 8 niveles.}
    \begin{tabular}{c|c}
        \hline
        \textbf{Intervalo de entrada} & \textbf{Nivel de cuantización} \\
        \hline
        $(-\infty, -3.5)$ & $-4$ \\
        $[-3.5, -2.5)$ & $-3$ \\
        $[-2.5, -1.5)$ & $-2$ \\
        $[-1.5, -0.5)$ & $-1$ \\
        $[-0.5, 0.5)$ & $0$ \\
        $[0.5, 1.5)$ & $1$ \\
        $[1.5, 2.5)$ & $2$ \\
        $[2.5, \infty)$ & $3$ \\
        \hline
    \end{tabular}
    \label{tab:quantization_example}
\end{table}

La tabla \ref{tab:quantization_example} define una función $q(x)$ que, por ejemplo, para un valor de entrada $x = 2.3$, devuelve un valor de salida $y = q(2.3) = 2$, ya que $2.3$ se encuentra en el intervalo $[1.5, 2.5)$, el cual está asociado al nivel de cuantización $2$.

En este caso los umbrales de cuantización quedan definidos por la \refeq{quantization_thresholds_example}.
\begin{equation}
    \label{eq:quantization_thresholds_example}
    \mathbf{t} = \{t_0, t_1, \ldots\} = \{-3.5, -2.5, -1.5, -0.5, 0.5, 1.5, 2.5\}.
\end{equation}

Y los niveles de cuantización por la \refeq{quantization_levels_example}.
\begin{equation}
    \label{eq:quantization_levels_example}
    \mathbf{l} = \{l_0, l_1, \ldots\} = \{-4, -3, -2, -1, 0, 1, 2, 3\}.
\end{equation}

Este tipo de cuantización, donde tanto los niveles como los umbrales de cuantización son valores equiespaciados, se conoce como cuantización uniforme \cite{sayood2017introduction}, la definición matemática formal de este tipo de cuantización se presenta en la sección \ref{sec:uniform}.

La cantidad de niveles $n_l$ en el caso de la cuantización uniforme depende del número de bits $b$ utilizados para representar cada valor cuantizado, y se calcula como $n_{l} = 2^b$. Por ejemplo, si se utilizan 8 bits para la cuantización, se tendrán $n_{l} = 2^8 = 256$ niveles de cuantización.

Otra función de cuantización utilizada normalmente es la cuantización no uniforme \cite{sayood2017introduction}. En donde los niveles y umbrales de cuantización no están equiespaciados y la cantidad de niveles no esta limitada a potencias de 2. Este tipo de cuantización puede ser beneficiosa en situaciones donde la distribución de los datos no es uniforme, permitiendo una mejor representación de ciertos rangos de valores. La misma será detallada en la sección \ref{sec:non_uniform}.

En la \reffig{uniform_vs_non_uniform} se muestra un ejemplo de ambos tipos de cuantización para un cuantizador de 8 niveles.

\TODO{Reemplazar figura con tikz}
\defaultfigure{cuantizacion/uniform_vs_non_uniform_3bits.png}{Cuantización uniforme y no uniforme.}{uniform_vs_non_uniform}

Otra característica importante es que la cuantización puede ser signada o no signada. En la cuantización signada, el rango de valores se divide en partes iguales alrededor de cero, mientras que en la cuantización no signada, los niveles de cuantización son siempre positivos, en donde se mejora la precisión para el rango positivo de valores. En la \reffig{sym_vs_asym} se muestra un ejemplo de ambos tipos de cuantización para un cuantizador uniforme de 3 bits.

\TODO{Reemplazar figura con tikz - sign/unsign}
\defaultfigure{cuantizacion/sym_vs_asym_3bits.png}{Cuantización simétrica y asimétrica.}{sym_vs_asym}

Toda función de cuantización introduce un error de cuantización, que es la diferencia entre el valor original y su representación cuantizada. Este error se busca minimizar para preservar la precisión de la representación original. En el caso, por ejemplo, de una cuantización uniforme de 8 niveles, el error de cuantización máximo es la mitad del tamaño del intervalo entre niveles consecutivos \cite{sayood2017introduction}. En aplicaciones como el procesamiento de señales, este error puede manifestarse como ruido, conocido como ruido de cuantización.

A diferencia de otras aplicaciones donde una métrica crítica es el error de cuantización, en la cuantización de redes neuronales el objetivo principal es minimizar el impacto de este error en la precisión del modelo, que si bien está relacionado con el mismo, no es equivalente. En este caso, el objetivo no es mantener la representación exacta de los pesos, si no una versión equivalente de la distribución de los pesos a efectos de la precisión del modelo. Analogamente a las compresiones de audio donde lo importante es mantener la calidad percibida y no la representación exacta de la señal original, con la ventaja de que en este caso se cuenta con un valor cuantitativo para medir la calidad de la compresión: la precisión del modelo.

Un aspecto importante, pero que no esta directamente definido por la función de cuantización es la representación numérica utilizada para almacenar los valores cuantizados. Generalmente se utiliza la notación INTX para referirse a representación en un formato de enteros (\textit{integer} en inglés) de X bits. Y la notación FPY para referirse a la representación en punto flotante (\textit{floating point} en inglés) de Y bits.

En el contexto de la cuantización de redes neuronales, como fue explicado en la sección \ref{sec:entrenamiento_redes_neuronales}, es necesaria una alta precisión numérica durante el entrenamiento de las mismas, por lo que este proceso se realiza generalmente en punto flotante de 32 bits (FP32). Luego, al aplicar cuantizaciones se suelen utilizar representaciones enteras, equivalentes a formatos de punto fijo en hardware (INT8, INT4).

En particular, la cuantización de redes neuronales consiste en representar los parámetros de la red con menor precisión numérica, mediante la conversión de sus parámetros $\theta$ representados originalmente en FP32 a formatos de menor precisión, como INT8 o INT4.

La principal razón para utilizar representaciones en punto fijo es la eficiencia computacional que ofrecen, ya que las operaciones en punto fijo son extremadamente eficientes en términos de velocidad y consumo energético en comparación con las operaciones en punto flotante. Esto se debe a que las operaciones en punto fijo requieren menos recursos computacionales y son más rápidas de ejecutar en hardware especializado, como unidades de procesamiento de tensores (TPU), unidades de procesamiento gráfico (GPU) o FPGA optimizadas para cálculos en punto fijo en donde generalmente se ejecutan modelos de aprendizaje automático. La desventaja que suele presentar la cuantización en punto fijo es la determinación de un rango adecuado para representar los valores, ya que un rango inadecuado puede llevar a saturaciones o pérdidas de precisión significativas \cite{gholami2021surveyquantizationmethodsefficient}.

A continuación, en la sección \ref{sec:funciones_de_cuantizacion}, se presentan las funciones de cuantización relevantes para la cuantización de redes neuronales, seguido por la descripción de las estrategias utilizadas para aplicar la cuantización en redes neuronales de una manera efectiva en la sección \ref{sec:estrategias_de_cuantizacion}.

\subsection{Funciones de cuantización}
\label{sec:funciones_de_cuantizacion}

\subsubsection{Cuantizador uniforme}
\label{sec:uniform}

El cuantizador uniforme consta de una serie de $n(b) = 2^b$ niveles equiespaciados dentro del rango $r(\alpha)$, donde $b$ es el número de bits empleados en la cuantización y $\alpha$ es un parámetro que define dicho rango.

El rango del cuantizador depende del tipo de cuantización:

\begin{itemize}
    \item En el caso signado, el rango es $r(\alpha) = 2\alpha$, abarcando el intervalo $[-\alpha, \alpha].$
    \item En el caso no signado, el rango es $r(\alpha) = \alpha$, abarcando el intervalo $[0, \alpha].$
\end{itemize}

De forma general, el intervalo del rango se define como \eqref{eq:range_interval}.

\begin{equation}
    \label{eq:range_interval}
    \mathbb{I} = [m, M]
\end{equation}


donde $m$ representa el valor mínimo y $M$ el valor máximo del intervalo del rango.

El paso de cuantización, denotado por $\Delta(\alpha, b)$, es la distancia entre niveles consecutivos y se define \eqref{eq:quantization_step}.

\begin{equation}
    \label{eq:quantization_step}
    \Delta(\alpha, b) = \frac{r(\alpha)}{n(b)}
\end{equation}

Los niveles de cuantización se expresan como \eqref{eq:uniform_quantizer_level_value}, donde $i$ representa el índice del nivel, que varía entre $0$ y $n(b)-1$.

\begin{equation}
    \label{eq:uniform_quantizer_level_value}
    l_i\left(\alpha; b\right) = m + i \cdot \Delta\left(\alpha, b\right) \text{, con } i = 0, 1, \ldots, n\left(b\right) - 1
\end{equation}

La función de cuantización uniforme se define \eqref{eq:uniform_quantizer} y su representación gráfica se muestra en la \reffig{uniform_signed}.

\begin{equation}
    \label{eq:uniform_quantizer}
    q_u(x; \alpha, b) =
    \begin{cases}
        l_{0} & \text{si $x \leq m$} \\
        \Delta \cdot \lfloor \frac{x}{\Delta} \rfloor & \text{si } m < x < M \\
        l_{n\left(b\right) - 1}\left(\alpha\right) & \text{si } x \geq M
    \end{cases}
\end{equation}

\defaultfigure{quantizers/uniform/q.png}{Función de cuantización uniforme signada.}{uniform_signed}

El cuantizador uniforme resulta sencillo de implementar tanto en hardware como en software debido a la distribución equiespaciada de sus niveles. Presenta eficiencia computacional, pues la cuantización se realiza mediante operaciones aritméticas simples, como divisiones y redondeos.

No obstante, existen escenarios en los que la cuantización uniforme no resulta adecuada.

Una aplicación donde los datos presentan una distribución no uniforme, esta puede no captar la variabilidad de los datos, lo que conduce a una pérdida significativa de información y a niveles en desuso. Si la aplicación exige alta precisión en determinadas regiones del rango, la cuantización uniforme puede no ofrecer la resolución necesaria en esas zonas críticas, este caso aparece con frecuencia en la distribución de los pesos de redes neuronales.

La cuantización uniforme resulta ineficiente al combinar rangos amplios con requisitos de alta precisión, puesto que la misma representación debe cubrir valores muy pequeños y muy grandes. Este enfoque puede ser viable con aritmética de punto flotante, no obstante, para aritmética entera o de punto fijo no es eficiente.

\subsubsection{Cuantizador no uniforme}
\label{sec:non_uniform}

El cuantizador no uniforme comprende una serie de $n$ niveles $\mathbf{l}$ cuyos límites están determinados por los umbrales $\mathbf{t}$. En este esquema no existen restricciones sobre la posición de umbrales y niveles, ni la noción de bits.

La función de cuantización no uniforme se define como \eqref{eq:non_uniform_quantizer} y su representación gráfica se muestra en la \reffig{non_uniform_signed_q}.

\begin{equation}
    \label{eq:non_uniform_quantizer}
    q_{nu}\left(x; \mathbf{t}, \mathbf{l}\right) =
    \begin{cases}
        l_{0} & \text{si $x \leq m$} \\
        l_{i} & \text{si } t_i < x < t_{i+1}, i = 0, 1, \ldots, n-2 \\
        l_{n-1} & \text{si } x \geq M
    \end{cases}
\end{equation}

donde $m$ y $M$ representan los límites del rango, es decir, $m = t_0$ y $M = t_{n-1}$.

\defaultfigure{quantizers/non_uniform/q.png}{Función de cuantización no uniforme signada.}{non_uniform_signed_q}

El cuantizador no uniforme es más versátil que el uniforme, pues permite adaptar los niveles a la distribución de los datos o a las necesidades del problema, lo que mejora la precisión y aumenta el nivel de compresión. No obstante, su implementación resulta más compleja, especialmente en hardware de punto fijo, debido a la incompatibilidad de algunos valores con la aritmética de punto fijo \cite{electronics13101923}.

% \TODO{Ver si hay que desarrollar xq es malo para punto fijo}

\subsubsection{Cuantizador flexible}
\label{sec:flex}

El cuantizador flexible pretende combinar la eficiencia de implementación del cuantizador uniforme con la capacidad del cuantizador no uniforme para adaptarse a distintas distribuciones de datos. Este método cuantiza uniformemente los niveles de un cuantizador no uniforme \cite{electronics13101923}.

El cuantizador flexible consta de un conjunto de $n$ niveles, con valores $\mathbf{l}$ cuyos límites están determinados por los umbrales $\mathbf{t}$. Como los casos anteriores, el rango se fija mediante el parámetro $\alpha$ y el rango del cuantizador depende del tipo de cuantización:

\begin{itemize}
    \item En el caso signado, el rango es $r(\alpha) = 2\alpha$, abarcando el intervalo $[-\alpha, \alpha].$
    \item En el caso no signado, el rango es $r(\alpha) = \alpha$, abarcando el intervalo $[0, \alpha].$
\end{itemize}

Para optimizar la implementación en la biblioteca, se definen $n$ niveles y $n+1$ umbrales. Esta elección puede resultar poco intuitiva, pero facilita la implementación posterior. En esta definición, los umbrales $t_0$ y $t_n$ son fijos y coinciden con los límites del rango: $t_0 = m$ y $t_n = M$.

La función de cuantización flexible se define \eqref{eq:flexible_quantizer}

\begin{equation}
    \label{eq:flexible_quantizer}
    q_{f}\left(x; \alpha, \mathbf{t}, \mathbf{l}\right) =
    \begin{cases}
        q_{u}\left(l_{0}; \alpha\right)& \text{si $x \leq m$} \\
        q_{u}\left(l_{i}; \alpha\right) & \text{si } t_i < x < t_{i+1}, i = 0, 1, \ldots, n-2 \\
        q_{u}\left(l_{n-1}; \alpha\right) & \text{si } x \geq M
    \end{cases}
\end{equation}

    donde $q_{u}$ es la función de cuantización uniforme definida \eqref{eq:uniform_quantizer}. El gráfico de la función de cuantización flexible se muestra en la \reffig{flexible_signed_q}.

\defaultfigure{quantizers/flex/q_q.png}{Función de cuantización flexible signada.}{flexible_signed_q}

La \reffig{flexible_signed_q} muestra, en amarillo, los niveles previos a la aplicación del cuantizador uniforme y, en azul, los niveles finales tras aplicar dicho cuantizador.

La principal diferencia respecto al uniforme es que los niveles $n$ no dependen de la cantidad de bits, tanto $\mathbf{l}$ como $\mathbf{t}$ son parámetros del cuantizador. En otras palabras, se desacoplan los niveles de cuantización de la cantidad de bits.

Este cuantizador aproxima la cuantización no uniforme mediante una función cuyos valores de salida son compatibles con la aritmética de punto fijo, por ello, permite una implementación eficiente en hardware en conjunto con la flexibilidad para adaptarse a distintas distribuciones de datos.


\subsection{Estrategias de cuantización}
\label{sec:estrategias_de_cuantizacion}
La cuantización reduce el tamaño del modelo, la memoria requerida y acelera la inferencia. No obstante, la reducción de la precisión numérica introduce un error de cuantización que puede afectar el rendimiento, desplazando el modelo del punto de convergencia obtenido en el entrenamiento previo a la cuantización. Por lo tanto, es esencial evaluar el impacto de la cuantización en las métricas y hallar un compromiso adecuado entre compresión y precisión.

Existen dos estrategias principales para aplicar la cuantización en redes neuronales: la cuantización post-entrenamiento (PTQ) y la cuantización durante el entrenamiento (QAT).

El objetivo de ambas estrategias es minimizar el impacto negativo de la cuantización en la precisión del modelo mediante el ajuste de los parámetros de cuantización y la adaptación del modelo a la representación cuantizada.

\subsubsection{Estrategia PTQ}
PTQ es la estrategia mayormente utilizada debida a su simplicidad y rapidez de implementación, ya que pueden determinarse los parámetros de cuantización a partir de una red ya entrenada, sin necesidad de reentrenamiento o ajuste fino \textit{fine-tuning} del modelo, y por ende no necesita un conjunto de datos de entrenamiento etiquetados para su aplicación \cite{gholami2021surveyquantizationmethodsefficient}.

Sin embargo, en modelos complejos o con cuantizaciones agresivas (por ejemplo, INT4), PTQ puede ocasionar pérdidas de precisión significativas debido a la incapacidad del modelo para adaptarse a la representación cuantizada \cite{dai2024trainablefixedpointquantizationdeep}.

Para seleccionar los parámetros de cuantización adecuados, se analiza la distribución de los pesos y activaciones de la red de la red. En el caso de los pesos, esto se realiza directamente sobre los valores de los pesos entrenados. Para analizar las activaciones, se utiliza un conjunto de datos representativo de las entradas para realizar una inferencia y recopilar estadísticas sobre las activaciones.

El proceso consiste en, a partir de una red entrenada, realizar los siguientes pasos:

\begin{enumerate}
    \item Calibración: uso de un conjunto de datos representativo para analizar la distribución de las activaciones de la red.
    \item Extracción de parámetros de cuantización: determinación de parámetros tales como niveles y umbrales a partir de las distribuciones de pesos y activaciones.
    \item Cuantización: aplicación de la función de cuantización a los parámetros de la red utilizando los parámetros obtenidos.
    \item Evaluación: evaluación del modelo cuantizado en un conjunto de prueba para medir el impacto sobre la precisión.
\end{enumerate}

Se ilustra el proceso de PTQ en la \reffig{ptq_process}.

\defaultdiagram{cuantizacion/ptq_process}{Cuantización post-entrenamiento (PTQ)}{ptq_process}{0.8}

La definición de los parámetros de cuantización es crítica, puesto que una calibración no representativa puede ocasionar pérdidas de precisión significativas.
\TODO{Mostrar distribucion en una figura}

\subsubsection{Estrategia QAT}
\label{sec:qat}
QAT consiste en entrenar un modelo considerando los efectos de la cuantización durante el proceso de entrenamiento \cite{jacob2017quantizationtrainingneuralnetworks}, al incorporarlos durante el entrenamiento, se permite que el modelo se adapte a la representación cuantizada, y generalmente produce modelos cuantizados de mayor precisión en comparación con PTQ, especialmente en cuantizaciones agresivas \cite{dai2024trainablefixedpointquantizationdeep}.

Como se explicó, en la sección \ref{sec:entrenamiento_redes_neuronales}, el proceso de entrenamiento de una red neuronal requiere de una alta precisión numérica, y suele realizarse en punto flotante, esta resolución numerica es particularmente crítica para el \textit{backward pass}, que depende de cálculos precisos para actualizar los pesos de manera efectiva y asegurar la convergencia del entrenamiento.

Las técnicas utilizadas para incluir los efectos de la cuantización durante el entrenamiento, se basan en emular los efectos de la cuantización durante la inferencia (o \textit{forward pass}), y luego realizar el \textit{backward pass} en representación de punto flotante.

Una técnica para emular estos efectos consiste en realizar una operación de cuantización sobre los tensores durante el \textit{forward pass}, y luego revertir esta cuantización \textit{dequantization} para el \textit{backward pass} \cite{dai2024trainablefixedpointquantizationdeep}. Este proceso se ilustra, para un tensor, en la \reffig{quant_dequant}.

\defaultdiagram{cuantizacion/quant_dequant}{Cuantización y descuantización}{quant_dequant}{0.9}

Otra técnica muy utilizada es \textit{fake-quantization} \cite{dai2024trainablefixedpointquantizationdeep}, que consiste en aplicar la función de cuantización durante el \textit{forward pass}, pero mantiene la representación numerica de los tensores en punto flotante. De esta manera, se utilizan valores numericos ajustados a los efectos de la cuantización, pero sin el costo computacional de realizar multiples operaciones de cuantización y descuantización. En este trabajo se utiliza la técnica de \textit{fake-quantization} para emular la cuantización durante el entrenamiento. Con esta técnica, no es necesaria la operación de descuantización, como se muestra en la \reffig{fake_quantization}, pues la representación numérica permanece en punto flotante.

\defaultdiagram{cuantizacion/fake_quantization}{Emulación de la cuantización (\textit{fake-quantization})}{fake_quantization}{0.9}
%  Sin embargo, incrementa el coste computacional y la duración del entrenamiento, ya que requiere de un reentrenamiento o ajuste fino (ajuste fino (fine-tuning)).
% \defaultdiagram{cuantizacion/qat_process}{Cuantización durante el entrenamiento (QAT)}{qat_process}{0.8}

QAT consiste en modificar el proceso de entrenamiento tradicional \cite{jacob2017quantizationtrainingneuralnetworks}, explicado en la sección \ref{sec:entrenamiento_redes_neuronales} y que se muestra en la \reffig{training}. QAT incorpora los efectos de la cuantización para forzar la adaptación del modelo a las limitaciones de la cuantización mediante el uso de \textit{fake-quantization} durante el \textit{forward pass}, para obtener una inferencia con los efectos de la cuantización, y luego, utilizarla para calcular la pérdida y retropropagar los gradientes, pero este proceso se mantiene en punto flotante para preservar la precisión del gradiente, como se muestra en la \reffig{qat}.

\defaultdiagram{cuantizacion/training}{Entrenamiento simple}{training}{1}
\defaultdiagram{cuantizacion/training_qat}{Entrenamiento con cuantización}{qat}{1}

En la \reffig{qat_algorithm} aparece la representación del término $\nabla_q$, correspondiente al gradiente de la función de cuantización. Sin embargo, las funciones de cuantización suelen tener puntos no diferenciables y derivada nula en el resto de su dominio (véase la sección \ref{sec:funciones_de_cuantizacion}), por lo que no es posible calcular el gradiente de forma directa mediante retropropagación pues esto impediría la propagación del gradiente e inhibiría la actualización de los pesos.

Como solución, se aproximan las funciones de cuantización $q$ por funciones diferenciables $\hat{q}$. La aproximación más simple y extendida es el \textit{Straight-Through Estimator} (STE) [CITA REQUERIDA], que reemplaza la derivada nula por una aproximación no nula en la región activa. Con esta aproximación, el gradiente se propaga a través de la cuantización como si fuera la identidad en la región activa. En la \reffig{q_ste} se muestra la función de cuantización uniforme (en verde) y su aproximación mediante STE (en rojo).

\defaultdiagram{cuantizacion/uniform_quantizer_with_ste}{Función de cuantización uniforme y su aproximación mediante STE}{q_ste}{1}

En este caso, se tiene la derivada de la función aproximada, como se muestra en la \refeq{ste_idea}.

\begin{equation}
    \label{eq:ste_idea}
    \frac{\partial q}{\partial x} \approx
    \frac{\partial \hat{q}}{\partial x} =
    \begin{cases}
        1 & x \in [-4, 4] \\
        0 & x \notin [-4, 4]
    \end{cases}
\end{equation}

En la \reffig{qat_algorithm} se muestra un diagrama detallado de un ejemplo del proceso de entrenamiento con QAT.

\defaultdiagram{cuantizacion/training_qat_example}{Ejemplo de entrenamiento con QAT}{qat_algorithm}{1}

El proceso de entrenamiento con QAT consiste en realizar el re-entrenamiento del modelo, mediante el entrenamiento con efectos de cuantización, este proceso de ilustra en la \reffig{qat_process}.

\defaultdiagram{cuantizacion/qat_process}{Proceso de cuantización durante el entrenamiento (QAT)}{qat_process}{0.8}
